%% ============================================================
\chapter{Nerode's Theorem --- Solutions}
%% ============================================================

\begin{enumerate}[{2.}1.]

%%------------------------------------------------------------
\item \textbf{Show $\Psi$ is not a right congruence.}

$x\,\Psi\,y \Leftrightarrow |x|-|y|$ is odd. Take $x=\mathbf{a}$, $y=\mathbf{b}$ (so $|x|-|y|=0$, not odd, thus $x\not\Psi\,y$). Also check: is $\Psi$ even an equivalence relation? $|x|-|x|=0$ (even), so $x\not\Psi\,x$ (not reflexive). Hence $\Psi$ is not an equivalence relation and a fortiori not a right congruence.

%%------------------------------------------------------------
\item \textbf{$Q$ defined by $|x|_\mathbf{a}-|y|_\mathbf{a}\equiv 0\pmod{3}$.}

\begin{enumerate}
\item \textit{Equivalence}: Reflexive: $|x|_\mathbf{a}-|x|_\mathbf{a}=0\equiv 0$. Symmetric: if $3\mid(|x|_\mathbf{a}-|y|_\mathbf{a})$ then $3\mid(|y|_\mathbf{a}-|x|_\mathbf{a})$. Transitive: if $3\mid(a-b)$ and $3\mid(b-c)$ then $3\mid(a-c)$.

\item \textit{Right congruence}: Suppose $x\,Q\,y$, i.e., $|x|_\mathbf{a}\equiv|y|_\mathbf{a}\pmod{3}$. For any $z\in\Sigma^*$:
$|xz|_\mathbf{a}=|x|_\mathbf{a}+|z|_\mathbf{a}$ and $|yz|_\mathbf{a}=|y|_\mathbf{a}+|z|_\mathbf{a}$. So $|xz|_\mathbf{a}-|yz|_\mathbf{a}=|x|_\mathbf{a}-|y|_\mathbf{a}\equiv 0\pmod{3}$, i.e., $xz\,Q\,yz$.
\end{enumerate}

%%------------------------------------------------------------
\item \textbf{All languages $L$ over $\{\mathbf{a},\mathbf{b},\mathbf{c}\}$ with $\mathrm{rk}(R_L)=1$.}

$\mathrm{rk}(R_L)=1$ means $R_L$ has exactly one equivalence class, namely all of $\Sigma^*$. This occurs iff $R_L = \Sigma^*\times\Sigma^*$, i.e., $x\,R_L\,y$ for all $x,y$. This happens iff $(\forall z)(xz\in L\Leftrightarrow yz\in L)$ for all $x,y,z$, which holds iff either $L=\emptyset$ or $L=\Sigma^*$.

\textit{Justification}: If $L\neq\emptyset$ and $L\neq\Sigma^*$, pick $w\in L$ and $u\not in L$. Then $w\,R_L\,u$ fails (take $z=\lambda$), so $R_L$ has at least 2 classes.

%%------------------------------------------------------------
\item \textbf{Show $P$ (with 3 equivalence classes as given) is a right congruence.}

The classes are $[\lambda]_P=\{\mathbf{0},\mathbf{1}\}^0$, $[\mathbf{1}]_P=\{\mathbf{0},\mathbf{1}\}^1$, $[\mathbf{00}]_P=\bigcup_{n\geq 2}\{\mathbf{0},\mathbf{1}\}^n$.

We must check: if $x\,P\,y$ then $xa\,P\,ya$ for all $a\in\{0,1\}$.

\textit{Case $x,y\in[\lambda]_P$}: $|x|=|y|=0$, so $|xa|=|ya|=1$, thus $xa,ya\in[\mathbf{1}]_P$.

\textit{Case $x,y\in[\mathbf{1}]_P$}: $|x|=|y|=1$, so $|xa|=|ya|=2$, thus $xa,ya\in[\mathbf{00}]_P$.

\textit{Case $x,y\in[\mathbf{00}]_P$}: $|x|,|y|\geq 2$, so $|xa|,|ya|\geq 3\geq 2$, thus $xa,ya\in[\mathbf{00}]_P$.

In every case $xa\,P\,ya$, so $P$ is a right congruence. $\square$

%%------------------------------------------------------------
\item \textbf{All languages $L$ with $R_L=P$ (from Exercise 2.4).}

Recall $A_0=[\lambda]_P=\{\lambda\}$, $A_1=[\mathbf{1}]_P=\Sigma^1$, $A_2=[\mathbf{00}]_P=\Sigma^{\geq 2}$.
For $R_L=P$ we need $L$ to be a union of $P$-classes \emph{and} no two distinct $P$-classes to be merged by $R_L$.
Two classes $A_i, A_j$ are merged (identified) by $R_L$ iff for every $z\in\Sigma^*$: $A_i\cdot z \subseteq L \Leftrightarrow A_j\cdot z\subseteq L$ (more precisely, for representative elements $u\in A_i$, $v\in A_j$: $uz\in L\Leftrightarrow vz\in L$ for all $z$).

Check each of the 6 nonempty proper unions:
\begin{itemize}
\item $L=A_0=\{\lambda\}$: For any $z\neq\lambda$, both $\mathbf{0}z\not in L$ (length$\geq 2$) and $\mathbf{00}z\not in L$; for $z=\lambda$, $\mathbf{0}\not in L$ and $\mathbf{00}\not in L$. So $A_1$ and $A_2$ are \emph{not} distinguished: $R_L\neq P$.
\item $L=A_1=\Sigma^1$: $z=\lambda$: $\lambda\not in L$ but $\mathbf{0}\in L$ (distinguishes $A_0$ from $A_1$). $\mathbf{0}\in L$ but $\mathbf{00}\not in L$ (distinguishes $A_1$ from $A_2$). All three classes remain distinct: $R_L=P$. \checkmark
\item $L=A_2=\Sigma^{\geq 2}$: $z=\mathbf{0}$: $\lambda\mathbf{0}=\mathbf{0}\not in L$ but $\mathbf{0}\cdot\mathbf{0}=\mathbf{00}\in L$ (distinguishes $A_0$ from $A_1$). $z=\lambda$: $\mathbf{0}\not in L$ but $\mathbf{00}\in L$ (distinguishes $A_1$ from $A_2$). So $R_L=P$. \checkmark
\item $L=A_0\cup A_1=\Sigma^{\leq 1}$: $z=\mathbf{0}$: $\lambda\cdot\mathbf{0}=\mathbf{0}\in L$ but $\mathbf{0}\cdot\mathbf{0}=\mathbf{00}\not in L$ (distinguishes $A_0$ from $A_1$). $z=\lambda$: $\mathbf{0}\in L$ but $\mathbf{00}\not in L$ (distinguishes $A_1$ from $A_2$). So $R_L=P$. \checkmark
\item $L=A_0\cup A_2$: $z=\lambda$: $\lambda\in L$ but $\mathbf{0}\not in L$ (distinguishes $A_0,A_1$). $z=\lambda$: $\mathbf{0}\not in L$ but $\mathbf{00}\in L$ (distinguishes $A_1,A_2$). So $R_L=P$. \checkmark
\item $L=A_1\cup A_2=\Sigma^+$: For any $z$, $\mathbf{0}z\in\Sigma^+$ and $\mathbf{00}z\in\Sigma^+$, both in $L$. So $A_1$ and $A_2$ are not distinguished: $R_L\neq P$.
\end{itemize}

Therefore exactly \textbf{4} languages satisfy $R_L=P$: $A_1$, $A_2$, $A_0\cup A_1$, and $A_0\cup A_2$.

%%------------------------------------------------------------
\item \textbf{All languages $L$ with $R_L=Q$ (from Exercise 2.2).}

$Q$ has three equivalence classes: $[x]_Q$ depends on $|x|_\mathbf{a}\bmod 3$. So $Q$-classes are $C_0=\{x:|x|_\mathbf{a}\equiv 0\}$, $C_1=\{x:|x|_\mathbf{a}\equiv 1\}$, $C_2=\{x:|x|_\mathbf{a}\equiv 2\}$. Any nonempty proper union of these classes gives $R_L=Q$.

%%------------------------------------------------------------
\item \textbf{Analysis of $Q$ on $\{\mathbf{a},\mathbf{b}\}^*$.}

\begin{enumerate}
\item \textit{Right congruence, classes}: $\lambda Q \lambda$ by definition; for non-empty strings, $x\,Q\,y$ iff $(|x|,|y|$ have the same parity. Classes: $[\lambda]_Q$ (just $\lambda$) and $C_e=\{x\neq\lambda:|x|\text{ even}\}$ and $C_o=\{x:|x|\text{ odd}\}$. This is a right congruence because appending any $a$ flips parity of length for non-empty strings, consistently within each class.

\item $L=[\lambda]_Q\cup[\mathbf{aa}]_Q=\{\lambda\}\cup C_e$. Simple description: $L=\{x:|x|\text{ is even}\}$. $R_L$ has two classes: even-length strings and odd-length strings.

\item $A_Q$: 3 states ($[\lambda]_Q$, $C_e$, $C_o$). Final states for $L$: $\{[\lambda]_Q, C_e\}$.

\item $A_{R_L}$: 2 states (even, odd), final: even.

\item $A_Q$ with relabeled final states: Yes. $C_e\cup[\lambda]_Q$ vs.\ $C_o$ -- this collapses to $A_{R_L}$.

\item The eight languages are the $2^3=8$ unions of subsets of $\{[\lambda]_Q, C_e, C_o\}$. A language $K$ satisfies the criterion of part (e) iff $R_K = Q$, i.e., $A_Q$ (relabeled with new final states for $K$) is already minimal.

Checking all eight: $\emptyset$ and $\Sigma^*$ give $\mathrm{rk}(R_K)=1$; $\{\lambda\}$, all-even $(\{\lambda\}\cup C_e)$, $C_o$, and $\Sigma^+$ $(=C_e\cup C_o)$ give $\mathrm{rk}(R_K)=2$ (two of the three $Q$-classes merge in each case). Only two choices keep all three $Q$-classes distinct:

\begin{itemize}
\item $K = C_e$: $z=\lambda$ distinguishes $[\lambda]$ from $C_e$ (since $\lambda\not in K$ but $aa\in K$); $z=\lambda$ distinguishes $C_e$ from $C_o$ (since $aa\in K$ but $a\not in K$). So $R_K=Q$.
\item $K = \{\lambda\}\cup C_o$: $z=aa$ distinguishes $[\lambda]$ from $C_o$ ($\lambda\cdot aa=aa\not in K$ but $a\cdot aa=aaa\in K$); $z=\lambda$ distinguishes $C_o$ from $C_e$ ($a\in K$ but $aa\not in K$). So $R_K=Q$.
\end{itemize}

These are the only two of the eight languages satisfying part (e).
\end{enumerate}

%%------------------------------------------------------------
\item \textbf{$I$ = identity on $\{\mathbf{a}\}^*$.}

\begin{enumerate}
\item $x\,I\,y\Leftrightarrow x=y$. Right congruence: if $x=y$, then $xa=ya$ for any $a$. Yes.
\item Each equivalence class is a singleton: $[\mathbf{a}^n]_I=\{\mathbf{a}^n\}$.
\end{enumerate}

%%------------------------------------------------------------
\item \textbf{$L=\{\lambda\}\cup\{\mathbf{a}\}\cup\{\mathbf{aa}\}$, $I$ has infinite rank. Is $L$ FAD?}

$L$ is a finite set, hence FAD. Nerode's theorem says $L$ is FAD iff $R_L$ has finite rank, not iff $I$ (or any other right congruence) has finite rank. Here $R_L$ has rank equal to the number of $R_L$-classes, which is finite (since $L$ is finite and FAD). The fact that $I$ has infinite rank is irrelevant -- we need $I$ to \emph{refine} $R_L$ (which it does, trivially, since $I$ is the finest relation), but Nerode's theorem requires $R_L$ itself to have finite rank. $\square$

%%------------------------------------------------------------
\item \textbf{Machine $M$ with $\mathrm{rk}(R_{L(M)})\neq |S_M|$.}

Let $M$ be the two-state DFA: $\Sigma=\{\mathbf{a}\}$, $S=\{q_0,q_1\}$, $s_0=q_0$, $F=\{q_0,q_1\}$ (all states final), $\delta(q_0,\mathbf{a})=q_1$, $\delta(q_1,\mathbf{a})=q_0$. Then $L(M)=\Sigma^*$, so $R_{L(M)}$ has rank 1, but $|S_M|=2$.

%%------------------------------------------------------------
\item \textbf{Show $F_{R_L}$ is well defined.}

$F_{R_L}=\{[x]_{R_L}\mid x\in L\}$. We must show: if $[x]_{R_L}=[y]_{R_L}$, then $x\in L\Leftrightarrow y\in L$ (so each $R_L$-class is either entirely inside $L$ or entirely outside $L$, making the assignment of classes to $F_{R_L}$ unambiguous). Since $[x]_{R_L}=[y]_{R_L}$, we have $x\,R_L\,y$, meaning $(\forall z\in\Sigma^*)(xz\in L\Leftrightarrow yz\in L)$. Taking $z=\lambda$: $x\in L\Leftrightarrow y\in L$. $\square$

%%------------------------------------------------------------
\item \textbf{Show $\delta_{R_L}$ is well defined.}

$\delta_{R_L}([x]_{R_L}, a) = [xa]_{R_L}$. We must show: if $[x]_{R_L}=[y]_{R_L}$, then $[xa]_{R_L}=[ya]_{R_L}$. Since $x\,R_L\,y$, for any $z$: $xaz\in L\Leftrightarrow x(az)\in L\Leftrightarrow y(az)\in L\Leftrightarrow yaz\in L$. Thus $xa\,R_L\,ya$. $\square$

%%------------------------------------------------------------
\item \textbf{Prove $\DBAR_{R_L}([x]_{R_L},y)=[xy]_{R_L}$ by induction on $|y|$.}

\textit{Base} $|y|=0$: $\DBAR_{R_L}([x],\lambda)=[x]=[x\lambda]$. \checkmark

\textit{Step}: $y=wa$, $|w|=n$. $\DBAR_{R_L}([x],wa)=\delta_{R_L}(\DBAR_{R_L}([x],w),a)=\delta_{R_L}([xw],a)=[xwa]=[xy]$. $\square$

%%------------------------------------------------------------
\item \textbf{Find $R^P$ for automaton $P$ in Example 2.6, check $R^P=R_L$.}

$R^P$ is defined by $x\,R^P\,y\Leftrightarrow\DBAR(s_0,x)=\DBAR(s_0,y)$. Since the machine $P$ in Example 2.6 is the minimal machine for $L$, each state is reachable and distinguishable. The relation $R^P$ has the same classes as $R_L$ (each state corresponds to one $R_L$-class). Hence $R^P=R_L$.

%%------------------------------------------------------------
\item \textbf{Find $R^A$ for machines from Chapter 1 exercises.}

For each DFA $A$ built in Chapter 1: $R^A$ is defined by $x\,R^A\,y\Leftrightarrow\DBAR(s_0,x)=\DBAR(s_0,y)$ (same state reached). Compare to $R_{L(A)}$: by definition of $R^A$, it always refines $R_{L(A)}$ (since if same state, same future acceptance). For minimal machines they coincide; for non-minimal ones $R^A$ may be strictly finer.

%%------------------------------------------------------------
\item \textbf{Prove the pumping lemma statement by induction.}

Let $M$ be a DFA with states $S$, $|S|=n$. For a word $w=uvw'$ with $|uv|\leq n$, $|v|\geq 1$, and $uvw'\in L$. Let $s=\DBAR(s_0,u)$, so $\DBAR(s,v)=s$ (since the path through $v$ loops back to $s$).

\textit{Claim}: $(\forall i\in\NN)(\DBAR(s_0,uv^iw')=\DBAR(s_0,uvw')=f\in F)$.

\textit{Proof by induction on $i$}:

$i=0$: $\DBAR(s_0,uw')=\DBAR(\DBAR(s_0,u),w')=\DBAR(s,w')$. And $\DBAR(s_0,uvw')=\DBAR(\DBAR(s_0,uv),w')=\DBAR(s,w')$ (since $\DBAR(s,v)=s$). Equal. \checkmark

$i\to i+1$: By the inner induction, $\DBAR(s_0,uv^i)=s$ (base: $\DBAR(s_0,u)=s$; step: $\DBAR(s_0,uv^{i+1})=\DBAR(\DBAR(s_0,uv^i),v)=\DBAR(s,v)=s$). Therefore
\[
\DBAR(s_0,uv^{i+1}w')=\DBAR(\DBAR(s_0,uv^i),vw')=\DBAR(s,vw')=\DBAR(\DBAR(s,v),w')=\DBAR(s,w').
\]

So $\DBAR(s_0,uv^iw')=\DBAR(s,w')=f$ for all $i$. $\square$

%%------------------------------------------------------------
\item \textbf{Prove Theorem 2.1.}

Theorem 2.1 states: $L$ is FAD iff $R_L$ has finite rank.

($\Rightarrow$) If $L=L(M)$ for a DFA $M$ with $n$ states, define $\mu: \Sigma^*\to S$ by $\mu(x)=\DBAR(s_0,x)$. Then $x\,R^M\,y\Leftrightarrow\mu(x)=\mu(y)$, so $R^M$ has at most $n$ classes. Since $R^M$ refines $R_L$ and $R_L$ is coarser, $\mathrm{rk}(R_L)\leq\mathrm{rk}(R^M)\leq n<\infty$.

($\Leftarrow$) If $R_L$ has finite rank $k$, construct the minimal machine $A_{R_L}=\langle\Sigma,[{\cdot}]_{R_L},[\lambda]_{R_L},\delta_{R_L},F_{R_L}\rangle$ with $k$ states (shown well-defined in Exercises 2.11--2.12). Show $L(A_{R_L})=L$: $x\in L(A_{R_L})\Leftrightarrow\DBAR_{R_L}([\lambda],x)\in F_{R_L}\Leftrightarrow[x]\in F_{R_L}\Leftrightarrow x\in L$.

%%------------------------------------------------------------
\item \textbf{Find a language giving rise to $I$ (identity on $\{\mathbf{a}\}^*$); show it cannot be FAD.}

\begin{enumerate}
\item Take
\[
L=\{\mathbf a^{2^m}\mid m\ge0\}\subseteq\{\mathbf a\}^*.
\]
We claim $R_L=I$.
Let $i<j$ and $d=j-i\ge1$.
Choose $m$ so that $2^m>d$, and set $z=\mathbf a^{2^m-i}$.
Then
\[
\mathbf a^i z=\mathbf a^{2^m}\in L,
\]
but
\[
\mathbf a^j z=\mathbf a^{2^m+d}
\]
with
\[
2^m<2^m+d<2^{m+1},
\]
so $2^m+d$ is not a power of 2. Hence $\mathbf a^j z\not in L$.
Thus $\mathbf a^i\not R_L\,\mathbf a^j$ for all $i\ne j$, so every class is a singleton and $R_L=I$.

\item Such a language cannot be FAD: $R_L=I$ has infinite rank (one class per string in $\{\mathbf{a}\}^*$), so by Nerode's theorem $L$ is not FAD. $\square$
\end{enumerate}

%%------------------------------------------------------------
\item \textbf{Prove Theorem 2.4 from Theorem 2.3.}

Theorem 2.3: If $L$ is FAD then $R_L$ has finite rank. Theorem 2.4 (Nerode): $L$ is FAD $\Leftrightarrow$ $R_L$ has finite rank.

We need the converse: finite rank $\Rightarrow$ FAD. Given $R_L$ has $k$ classes, the machine $A_{R_L}$ has $k$ states and accepts $L$ (by Exercises 2.11--2.13). Hence $L$ is FAD.

%%------------------------------------------------------------
\item \textbf{Prove Theorem 2.5.}

Theorem 2.5: $L$ is FAD $\Leftrightarrow$ there exists a right congruence $R$ of finite rank such that $L$ is a union of $R$-classes.

($\Rightarrow$) Take $R=R^M$ for a DFA $M$ accepting $L$: finite rank (at most $|S|$), and $L=\bigcup_{s\in F}[s]_{R^M}$.

($\Leftarrow$) Given such $R$, build $A_R$ with states = $R$-classes, $F_R$ = classes $\subseteq L$. Then $L=L(A_R)$, so $L$ is FAD.

%%------------------------------------------------------------
\item \textbf{Prove Theorem 2.6.}

Theorem 2.6 states:
\[
(\exists n)(\forall x\not in L,\ |x|\ge n)(\exists u,v,w)\bigl[x=uvw,\ |uv|\le n,\ |v|\ge1,\ (\forall i\ge0)\ uv^iw\not in L\bigr].
\]

Assume $L$ is FAD, accepted by DFA $M=\langle\Sigma,S,s_0,\delta,F\rangle$, with $|S|=n$.
Take any $x\not in L$ with $|x|\ge n$.
Look at the first $n+1$ states on the run of $x$:
\[
s_0,\ \DBAR(s_0,a_1),\ \DBAR(s_0,a_1a_2),\ \ldots,\ \DBAR(s_0,a_1\cdots a_n).
\]
By pigeonhole, two are equal. So for some decomposition $x=uvw$ with $|uv|\le n$, $|v|\ge1$,
\[
\DBAR(s_0,u)=\DBAR(s_0,uv)=q.
\]
Hence $v$ is a loop at $q$, so for every $i\ge0$,
\[
\DBAR(s_0,uv^iw)=\DBAR(q,w)=\DBAR(s_0,uvw)=\DBAR(s_0,x).
\]
Since $x\not in L$, $\DBAR(s_0,x)\not in F$. Therefore $\DBAR(s_0,uv^iw)\not in F$, i.e.,
$uv^iw\not in L$ for all $i\ge0$. This is exactly Theorem 2.6. $\square$

%%------------------------------------------------------------
\item \textbf{Prove Theorem 2.7.}

Theorem 2.7 states: if $M$ is an $n$-state DFA accepting $L$, then every
$x=a_1a_2\cdots a_m\in L$ with $m\ge n$ has a subsequence
$a_{i_1}a_{i_2}\cdots a_{i_j}\in L$ with $j<n$.

Let $x\in L$, $|x|=m\ge n$.
By Theorem 2.3 (or Exercise 2.16), we can write
\[
x=u_1v_1w_1,\quad |u_1v_1|\le n,\ |v_1|\ge1,\ \text{and }u_1v_1^iw_1\in L\ \forall i\ge0.
\]
In particular, with $i=0$, $x_1=u_1w_1\in L$ and $|x_1|<|x|$.

If $|x_1|<n$, stop. Otherwise apply the same argument to $x_1$.
Each step removes at least one symbol and keeps membership in $L$, so after finitely many steps we obtain
\[
y\in L,\quad |y|<n,
\]
and $y$ is obtained from $x$ by deleting some blocks.
Therefore $y$ is a subsequence of $x$:
\[
y=a_{i_1}a_{i_2}\cdots a_{i_j}
\]
for an increasing index sequence $i_1<i_2<\cdots<i_j$, with $j=|y|<n$.
This is exactly Theorem 2.7. $\square$

%%------------------------------------------------------------
\item \textbf{Show $L=\{x\mid|x|_\mathbf{a}<|x|_\mathbf{b}\}$ is not FAD.}

For any $n$, take $w=\mathbf{a}^n\mathbf{b}^{n+1}\in L$, $|w|=2n+1\geq n$. Any decomposition $w=uvw'$ with $|uv|\leq n$ and $|v|\geq 1$ has $v=\mathbf{a}^k$ for some $k\geq 1$ (since $|uv|\leq n$). Then $uv^0w'=uw'=\mathbf{a}^{n-k}\mathbf{b}^{n+1}$ which has $n-k$ $\mathbf{a}$s and $n+1$ $\mathbf{b}$s. Since $n-k<n+1$, this is in $L$. And $uv^2w'=\mathbf{a}^{n+k}\mathbf{b}^{n+1}$: need $n+k < n+1$, i.e., $k<1$, impossible since $k\geq 1$. So $uv^2w'\not in L$. The pumping lemma fails, so $L$ is not FAD. $\square$

%%------------------------------------------------------------
\item \textbf{Show $G=\{x\mid|x|_\mathbf{a}\geq|x|_\mathbf{b}\}$ is not FAD.}

For any $n$, take $w=\mathbf{a}^n\mathbf{b}^n\in G$. Any $v=\mathbf{a}^k$ ($k\geq 1$, since $|uv|\leq n$). Then $uv^0w'=\mathbf{a}^{n-k}\mathbf{b}^n$: need $n-k\geq n$, i.e., $k\leq 0$. Contradiction. So pumping down takes us out of $G$. Not FAD. $\square$

%%------------------------------------------------------------
\item \textbf{Show $P=\{\mathbf{d}^p\mid p\text{ prime}\}$ is not FAD.}

Suppose $P$ is FAD with $n$ states. Take a prime $p>n$ and $w=\mathbf{d}^p\in P$. By pumping: $w=uvw'$ with $|uv|\leq n$, $|v|=k\geq 1$, so $v=\mathbf{d}^k$ and $u=\mathbf{d}^j$, $w'=\mathbf{d}^{p-j-k}$. Then $uv^iw'=\mathbf{d}^{j+ik+(p-j-k)}=\mathbf{d}^{p+(i-1)k}$ must be in $P$ for all $i\geq 0$. So $p+(i-1)k$ must be prime for all $i\geq 0$. Take $i=p+1$: $p+pk=p(1+k)$, which is composite (both factors $>1$ since $p\geq 2$ and $k\geq 1$). Contradiction. $\square$

%%------------------------------------------------------------
\item \textbf{Show $\Gamma=\{w\cdot\mathbf{2}\cdot w\mid w\in\{\mathbf{0},\mathbf{1}\}^*\}$ is not FAD.}

We use Theorem 2.6. Suppose $\Gamma$ were FAD, and let $n$ be the number
supplied by that theorem.

Take
\[
x=\mathbf{0}^{\,n!}\mathbf{2}\mathbf{0}^{\,2n!}.
\]
This string is \emph{not} in $\Gamma$, because the block to the left of
$\mathbf{2}$ has length $n!$ while the block to the right has length $2n!$, so
the two sides are unequal.

By Theorem 2.6 there should be a decomposition
\[
x=uvw,\qquad |uv|\le n,\qquad |v|\ge1,
\]
such that
\[
uv^iw\not in\Gamma\qquad\text{for every }i\ge0.
\]

Since the first $n$ symbols of $x$ are all $\mathbf{0}$'s, the substring $v$
must be of the form $\mathbf{0}^k$ for some $1\le k\le n$.
Because $k\le n$, the integer $k$ divides $n!$.
Choose
\[
i=\frac{n!}{k}+1.
\]
Then pumping $v$ changes only the block of $\mathbf{0}$'s to the left of
$\mathbf{2}$, and its new length is
\[
n!+(i-1)k=n!+\frac{n!}{k}k=2n!.
\]
Therefore
\[
uv^iw=\mathbf{0}^{\,2n!}\mathbf{2}\mathbf{0}^{\,2n!}\in\Gamma.
\]
This contradicts Theorem 2.6.

Hence $\Gamma$ is not FAD. $\square$

%%------------------------------------------------------------
\item \textbf{Show $\Psi=\{ww\mid w\in\{\mathbf{0},\mathbf{1}\}^*\}$ is not FAD.}

For each $n\ge 0$, let $x_n=\mathbf{0}^n\mathbf{1}$.
We show the $x_n$ are pairwise in distinct $R_\Psi$-classes.

Fix $i\neq j$ and take suffix $z_i=\mathbf{0}^i\mathbf{1}$.
Then
\[
x_i z_i=\mathbf{0}^i\mathbf{1}\mathbf{0}^i\mathbf{1}=ww\in\Psi
\]
with $w=\mathbf{0}^i\mathbf{1}$.

Now consider
\[
x_j z_i=\mathbf{0}^j\mathbf{1}\mathbf{0}^i\mathbf{1}.
\]
Suppose this were in $\Psi$, so $x_j z_i=tt$ for some $t$.
The string has exactly two $\mathbf{1}$s, so each half must contain exactly one $\mathbf{1}$.
Hence the position of the second $\mathbf{1}$ must be ``half-length'' after the first.
Its first $\mathbf{1}$ is at position $j+1$, second at $j+i+2$, so the half-length must be $i+1$.
But half-length is also $\frac{i+j+2}{2}$, giving $\frac{i+j+2}{2}=i+1$, hence $j=i$, contradiction.
Therefore $x_j z_i\not in\Psi$.

So $x_i\not R_\Psi x_j$ for $i\neq j$, giving infinitely many Nerode classes. Thus $\Psi$ is not FAD. $\square$

%%------------------------------------------------------------
\item \textbf{Show $K=\{w\mid w=w^r\}$ (palindromes) is not FAD.}

For any $n$, the strings $\mathbf{0},\mathbf{0}^2,\ldots,\mathbf{0}^n$ are in distinct $R_K$-classes: $\mathbf{0}^i\,R_K\,\mathbf{0}^j$ requires $(\forall z)(\mathbf{0}^iz\in K\Leftrightarrow\mathbf{0}^jz\in K)$. Take $z=\mathbf{0}^i$: $\mathbf{0}^{2i}$ is a palindrome (yes). $\mathbf{0}^{j+i}$: palindrome (yes, any string of one letter is a palindrome). So we need a subtler argument. Take $z=\mathbf{1}\mathbf{0}^i$: $\mathbf{0}^i\mathbf{1}\mathbf{0}^i\in K$ (palindrome), but $\mathbf{0}^j\mathbf{1}\mathbf{0}^i\in K$ iff $\mathbf{0}^j\mathbf{1}\mathbf{0}^i=\mathbf{0}^i\mathbf{1}\mathbf{0}^j$, i.e., $i=j$. So distinct $i\neq j$ give distinct classes. Infinite rank. Not FAD. $\square$

%%------------------------------------------------------------
\item \textbf{Show $\Phi=\{\mathbf{a}^j\mathbf{b}^k\mathbf{c}^m\mid j\geq 3,k=m\}$ is not FAD (second pumping lemma).}

We use Theorem 2.5. Suppose, for contradiction, that $\Phi$ is FAD. Then there is a
right congruence $R$ of finite rank such that $\Phi$ is a union of $R$-classes.

For each $i\in\NN$, let
\[
x_i=\mathbf{a}^3\mathbf{b}^i.
\]
We claim that the $x_i$ lie in distinct $R$-classes. If $i\neq j$ and $x_i\,R\,x_j$, then
by the right-congruence property,
\[
x_i\mathbf{c}^i\,R\,x_j\mathbf{c}^i.
\]
But
\[
x_i\mathbf{c}^i=\mathbf{a}^3\mathbf{b}^i\mathbf{c}^i\in\Phi,
\]
while
\[
x_j\mathbf{c}^i=\mathbf{a}^3\mathbf{b}^j\mathbf{c}^i\not in\Phi
\]
because $j\neq i$.
Since $\Phi$ is a union of $R$-classes, two $R$-equivalent strings cannot disagree
about membership in $\Phi$. Contradiction.

So $x_0,x_1,x_2,\ldots$ belong to infinitely many distinct $R$-classes, contradicting
the finite rank of $R$. Therefore $\Phi$ is not FAD. $\square$

\textit{Why first pumping lemma is hard}: Any long word in $\Phi$ starts with $\geq 3$ $\mathbf{a}$s. Pumping within the $\mathbf{a}$-block may still satisfy $j\geq 3$; pumping within $\mathbf{b}$ or $\mathbf{c}$ block destroys $k=m$.

%%------------------------------------------------------------
\item \textbf{Show $C=\{\mathbf{d}^q\mid q\text{ non-prime}\}$ is not FAD.}

We use Theorem 2.6. Suppose $C$ were FAD, and let $n$ be the number supplied by
that theorem.

Choose a prime $p>n$ and let
\[
x=\mathbf{d}^p.
\]
Since $p$ is prime, $x\not in C$. By Theorem 2.6 there should be a decomposition
\[
x=uvw,\qquad |uv|\le n,\qquad |v|\ge1,
\]
such that
\[
uv^iw\not in C\qquad\text{for every }i\ge0.
\]

Because $x$ is unary, $v=\mathbf{d}^k$ for some $k\ge1$. Then
\[
uv^iw=\mathbf{d}^{\,p+(i-1)k}.
\]
Now choose $i=p+1$. Its length is
\[
p+pk=p(1+k),
\]
which is composite. Hence
\[
uv^{p+1}w\in C,
\]
contradicting Theorem 2.6.

Therefore $C$ is not FAD. $\square$

%%------------------------------------------------------------
\item \textbf{$R_L$ with 3 classes: $\{\lambda\}$, odd-length, even-length$\setminus\{\lambda\}$.}

\begin{enumerate}
\item Why $L=\{x\mid|x|\text{ odd}\}$ can't give this $R_L$: Compute $R_{\{x:|x|\text{ odd}\}}$. We need $x\,R_L\,y\Leftrightarrow(\forall z)(|xz|\text{ odd}\Leftrightarrow|yz|\text{ odd})$, i.e., $|x|\equiv|y|\pmod 2$. So $R_L$ has only 2 classes (even-length, odd-length), not 3. $\lambda$ is not distinguished from other even-length strings.

\item Let
\[
A=\{\lambda\},\quad B=\{x:|x|\text{ odd}\},\quad C=\{x:|x|\text{ even and }x\neq\lambda\}.
\]
Checking all 8 unions of $A,B,C$, exactly two give this same 3-class relation:
\[
L=C \quad\text{and}\quad L=A\cup B.
\]
For each of these, $A,B,C$ are pairwise distinguishable:
\begin{itemize}
\item $A$ vs $C$: suffix $z=\lambda$.
\item $A$ vs $B$: suffix $z$ of odd length.
\item $B$ vs $C$: suffix $z=\lambda$.
\end{itemize}
All other unions merge at least one pair of classes.
\end{enumerate}

%%------------------------------------------------------------
\item \textbf{Describe $R_\Psi$ and draw machine for $\Psi=\{x\mid|x|_\mathbf{a}\text{ even, ends with }\mathbf{b}\}$.}

$R_\Psi$ classes: we need to determine when $x,y$ have the same future w.r.t.\ $\Psi$. The future of $x$ depends on: (1) parity of $|x|_\mathbf{a}$, and (2) whether $x$ ends with $\mathbf{b}$.

Thus $R_\Psi$ has four classes:
\begin{itemize}
\item $C_{e,b}=\{x\mid |x|_\mathbf{a}\text{ even and }x\text{ ends with }\mathbf{b}\}$.
\item $C_{e,\neg b}=\{x\mid |x|_\mathbf{a}\text{ even and }x\text{ does not end with }\mathbf{b}\}$.
\item $C_{o,b}=\{x\mid |x|_\mathbf{a}\text{ odd and }x\text{ ends with }\mathbf{b}\}$.
\item $C_{o,\neg b}=\{x\mid |x|_\mathbf{a}\text{ odd and }x\text{ does not end with }\mathbf{b}\}$.
\end{itemize}
Here $\lambda\in C_{e,\neg b}$ (even number of $\mathbf a$'s, and it does not end with $\mathbf b$).

So the minimal DFA has 4 states, start state $C_{e,\neg b}$, and unique final state $C_{e,b}$.

%%------------------------------------------------------------
\item \textbf{Show $\Xi=\{\lambda,\mathbf{a},\mathbf{a}^4,\mathbf{a}^9,\ldots\}$ (perfect square lengths) is not FAD.}

We use Theorem 2.3. Suppose $\Xi$ were FAD, and let $n$ be a pumping length for $\Xi$.
Take
\[
w=\mathbf{a}^{n^2}\in\Xi.
\]
For any decomposition
\[
w=uvw,\qquad |uv|\le n,\qquad |v|\ge1,
\]
the string $v$ must be of the form $\mathbf{a}^k$ with $1\le k\le n$.

Pump once upward:
\[
uv^2w=\mathbf{a}^{n^2+k}.
\]
Since $1\le k\le n$, we have
\[
n^2<n^2+k<n^2+(2n+1)=(n+1)^2.
\]
So $n^2+k$ is strictly between two consecutive squares, hence is not itself a square.
Therefore
\[
uv^2w\not in\Xi,
\]
contradicting Theorem 2.3.

Thus $\Xi$ is not FAD. $\square$

%%------------------------------------------------------------
\item \textbf{Show $\Phi=\{\mathbf{b},\mathbf{b}^2,\mathbf{b}^4,\mathbf{b}^8,\ldots\}$ (powers of 2) is not FAD.}

We use Theorem 2.3. Suppose $\Phi$ were FAD, and let $n$ be a pumping length for $\Phi$.
Choose $m$ so that $2^m>n$, and let
\[
w=\mathbf{b}^{2^m}\in\Phi.
\]
For any decomposition
\[
w=uvw,\qquad |uv|\le n,\qquad |v|\ge1,
\]
we have $v=\mathbf{b}^k$ for some $1\le k\le n<2^m$.

Pump once upward:
\[
uv^2w=\mathbf{b}^{2^m+k}.
\]
Because $0<k<2^m$,
\[
2^m<2^m+k<2^{m+1}.
\]
So $2^m+k$ lies strictly between two consecutive powers of $2$, and therefore is not
itself a power of $2$. Hence
\[
uv^2w\not in\Phi,
\]
contradicting Theorem 2.3.

Therefore $\Phi$ is not FAD. $\square$

%%------------------------------------------------------------
\item \textbf{Five-class $R_L$ over $\{\mathbf{a},\mathbf{b}\}$.}

\begin{enumerate}
\item Let
\[
C_0=\{\lambda\},\; C_1=\{\mathbf a\},\; C_2=\{\mathbf{aa}\},\; C_3=\{\mathbf a^k:k\ge 3\},\; C_4=\{x:x\text{ contains a }\mathbf b\}.
\]
If $L$ is exactly one class and still has $R_L$ equal to this 5-class partition, the only possibility is
\[
L=C_3=\{\mathbf a^3,\mathbf a^4,\ldots\}.
\]
Choosing any other single class merges some classes under $R_L$.

\item The other languages $L$ that still give this same $R_L$ are exactly:
\[
C_0\cup C_3,\quad C_1\cup C_3,\quad C_0\cup C_1\cup C_3,
\]
\[
C_2\cup C_4,\quad C_0\cup C_2\cup C_4,\quad C_1\cup C_2\cup C_4,\quad C_0\cup C_1\cup C_2\cup C_4.
\]
Together with part (a), there are 8 such languages.
\end{enumerate}

%%------------------------------------------------------------
\item \textbf{Find $R_\Omega$ for $\Omega=\{y\mid(y\text{ has exactly one }\mathbf{0})\vee(y\text{ has even \#1s})\}$.}

$xz\in\Omega$ iff $(|x|_\mathbf{0}+|z|_\mathbf{0}=1)$ or $(|x|_\mathbf{1}+|z|_\mathbf{1}\equiv 0\pmod 2)$.
The future of $x$ (i.e., which $z$ satisfy $xz\in\Omega$) depends only on $c_0=|x|_\mathbf{0}$ (capped at $2$) and $p=|x|_\mathbf{1}\bmod 2$.
Writing $(c_0,p)$ for these two pieces of information, the six cases yield the following future descriptions:

\begin{itemize}
\item $(0,0)$: $xz\in\Omega\Leftrightarrow |z|_\mathbf{0}=1$ or $|z|_\mathbf{1}$ even.
\item $(0,1)$: $xz\in\Omega\Leftrightarrow |z|_\mathbf{0}=1$ or $|z|_\mathbf{1}$ odd.
\item $(1,0)$: $xz\in\Omega\Leftrightarrow |z|_\mathbf{0}=0$ or $|z|_\mathbf{1}$ even.
\item $(1,1)$: $xz\in\Omega\Leftrightarrow |z|_\mathbf{0}=0$ or $|z|_\mathbf{1}$ odd.
\item $(2,0)$: $xz\in\Omega\Leftrightarrow |z|_\mathbf{1}$ even. \quad ($|xz|_\mathbf{0}\geq 2$ so the first disjunct can never hold.)
\item $(2,1)$: $xz\in\Omega\Leftrightarrow |z|_\mathbf{1}$ odd.
\end{itemize}

All six futures are distinct (for example, $(0,0)$ and $(2,0)$ differ on $z=\mathbf{01}$: the first accepts it since $|z|_\mathbf{0}=1$, the second rejects since $|z|_\mathbf{1}=1$ is odd; and so on for the other pairs). Hence $R_\Omega$ has exactly \textbf{6} equivalence classes:
\[
C_{c_0,p}=\{x\in\{\mathbf{0},\mathbf{1}\}^*\mid |x|_\mathbf{0}=c_0\text{ (or }\geq 2\text{ if }c_0=2)\text{ and }|x|_\mathbf{1}\equiv p\pmod 2\},
\]
for $(c_0,p)\in\{0,1,2\}\times\{0,1\}$.
The minimal DFA for $\Omega$ has 6 states. Final states are those in $\Omega$: $(1,0),(1,1),(0,0),(2,0)$ — i.e., strings with exactly one $\mathbf{0}$ (states $(1,\cdot)$) or with even \#$\mathbf{1}$s (states $(\cdot,0)$).

%%------------------------------------------------------------
\item \textbf{Which of $L_1,L_2,L_1\cap L_2,\sim L_2,L_1\cup L_2$ are FAD?}

$L_1=\{x\mid|x|_\mathbf{a}>|x|_\mathbf{b}\}$: \textbf{not FAD}. If it were FAD, let $n$ be a pumping length from Theorem 2.3. Take
\[
w=\mathbf{a}^{n+1}\mathbf{b}^n\in L_1.
\]
Any split $w=uvw'$ with $|uv|\le n$ and $|v|\ge1$ has $v=\mathbf{a}^k$ for some $k\ge1$.
Pumping down gives
\[
uw'=\mathbf{a}^{n+1-k}\mathbf{b}^n,
\]
and now $n+1-k\le n$, so this pumped string is not in $L_1$. Contradiction.

$L_2=\{x\mid|x|_\mathbf{a}<3\}$: \textbf{FAD}. A DFA only needs to remember whether it has seen $0$, $1$, $2$, or at least $3$ copies of $\mathbf a$.

$L_1\cap L_2$: \textbf{FAD}. In fact this language is finite:
\[
L_1\cap L_2=\{\mathbf a,\mathbf{aa},\mathbf{aab},\mathbf{aba},\mathbf{baa}\}.
\]
Every finite language is FAD.

$\sim L_2=\{x\mid|x|_\mathbf{a}\geq 3\}$: \textbf{FAD}. The same 4-state machine used for $L_2$ works here too; we simply make the ``at least 3 $\mathbf a$'s'' state final.

$L_1\cup L_2$: \textbf{not FAD}. Suppose it were FAD, and let $n$ be a pumping length. Consider
\[
w=\mathbf{a}^{n+4}\mathbf{b}^{n+3}\in L_1\cup L_2
\]
because it has more $\mathbf a$'s than $\mathbf b$'s.
Again any split $w=uvw'$ with $|uv|\le n$ and $|v|\ge1$ has $v=\mathbf{a}^k$ for some $k\ge1$.
Pumping down gives
\[
uw'=\mathbf{a}^{n+4-k}\mathbf{b}^{n+3}.
\]
Here $n+4-k\ge4$, so the pumped string is not in $L_2$; and also $n+4-k\le n+3$, so it is not in $L_1$.
Thus $uw'\not in L_1\cup L_2$, contradicting Theorem 2.3.

%%------------------------------------------------------------
\item \textbf{$R_m$ defined by $|x|-|y|$ is a multiple of $m$.}

\begin{enumerate}
\item \textit{Equivalence relation and right congruence}: Reflexive, symmetric, and transitive are immediate from divisibility by $m$. If $m\mid(|x|-|y|)$, then
\[
|xz|-|yz|=|x|-|y|,
\]
so $m\mid(|xz|-|yz|)$ as well. Thus $R_m$ is a right congruence.

\item $R_2\cap R_3 = R_6$: $x\,R_2\cap R_3\,y\Leftrightarrow 2\mid|x|-|y|$ and $3\mid|x|-|y|\Leftrightarrow 6\mid|x|-|y|\Leftrightarrow x\,R_6\,y$.
So every $R_6$-class is contained in an $R_2$-class (and also in an $R_3$-class); in particular, $R_6\subseteq R_2$.

\item In general: since both $R$ and $S$ are equivalence relations, $R\cap S$ is also an equivalence relation. And if $x\,R\cap S\,y$, then $x\,R\,y$ and $x\,S\,y$. For any $z$, right congruence of $R$ and $S$ gives $xz\,R\,yz$ and $xz\,S\,yz$, hence $xz\,(R\cap S)\,yz$. So $R\cap S$ is a right congruence. Also $R\cap S\subseteq R$ and $R\cap S\subseteq S$, so its classes refine the classes of both $R$ and $S$.

\item $R_2\cup R_3$: $\lambda\,R_2\,\mathbf{aa}$ and $\mathbf{a}\,R_3\,\mathbf{aaaa}$. But is $\lambda\,R_2\cup R_3\,\mathbf{a}$? $|\lambda|-|\mathbf{a}|=-1$, not divisible by 2 or 3. No. Is $\lambda\,R_2\cup R_3\,\mathbf{aaa}$? $0-3=-3$, divisible by 3. Yes. Now $\lambda\,R_2\cup R_3\,\mathbf{aaa}$ and $\mathbf{aaa}\,R_2\cup R_3\,\mathbf{aaaaa}$ (length diff=2). But is $\lambda\,R_2\cup R_3\,\mathbf{aaaaa}$? Length diff=5, not divisible by 2 or 3. No. So transitivity fails: not an equivalence relation.

\item If $R\cup S$ is an equivalence relation: for any $x,y$ with $x\,R\cup S\,y$, we have $x\,R\,y$ or $x\,S\,y$. Either way, for any $z$: $xz\,R\,yz$ or $xz\,S\,yz$ (by right congruence of $R$ or $S$), hence $xz\,R\cup S\,yz$.
\end{enumerate}

%%------------------------------------------------------------
\item \textbf{Example of two right congruences whose union is not a right congruence.}

Let $R=R_2$ and $S=R_3$. As shown above, $R_2\cup R_3$ is not even an equivalence relation, hence not a right congruence.

%%------------------------------------------------------------
\item \textbf{Analysis of $\Gamma=\{\lambda,\mathbf{a},\mathbf{ab},\mathbf{ba},\mathbf{bb},\mathbf{bbb}\}\cup\{x:|x|\geq 4\}$.}

\begin{enumerate}
\item $\mathbf{ab}\,R_\Gamma\,\mathbf{ba}$: For all $z$: $\mathbf{ab}z\in\Gamma\Leftrightarrow\mathbf{ba}z\in\Gamma$. If $|z|=0$: $\mathbf{ab},\mathbf{ba}\in\Gamma$ (both listed). \checkmark If $|z|=1$: $|\mathbf{ab}z|=3$. Since $\Gamma$ contains only $\mathbf{bbb}$ at length 3, and $\mathbf{ab}$ cannot be extended to $\mathbf{bbb}$ by appending one letter, $\mathbf{ab}z\not in\Gamma$; likewise $\mathbf{ba}z\not in\Gamma$ (since $\mathbf{baa},\mathbf{bab}\neq\mathbf{bbb}$). Both outside $\Gamma$. \checkmark If $|z|\geq 2$: $|\mathbf{ab}z|,|\mathbf{ba}z|\geq 4$, so both in $\Gamma$. \checkmark Hence $\mathbf{ab}\,R_\Gamma\,\mathbf{ba}$.

\item $\mathbf{ab}$ is not related to $\mathbf{bb}$: take $z=\mathbf{b}$. Then
\[
\mathbf{ab}z=\mathbf{abb}\not in\Gamma
\]
because among length-3 strings only $\mathbf{bbb}$ is in $\Gamma$.
But
\[
\mathbf{bb}z=\mathbf{bbb}\in\Gamma.
\]
So $\mathbf{ab}\not R_\Gamma\,\mathbf{bb}$. $\square$

\item The 8 classes are exactly:
\[
C_0=\{\lambda\},\ C_1=\{\mathbf a\},\ C_2=\{\mathbf b\},\ C_3=\{\mathbf{aa}\},\ C_4=\{\mathbf{bb}\},
\]
\[
C_5=\{\mathbf{ab},\mathbf{ba}\},\ C_6=\{x:|x|=3,\ x\neq \mathbf{bbb}\},\ C_7=\{\mathbf{bbb}\}\cup\{x:|x|\ge4\}.
\]
These are pairwise distinct by suitable suffixes (for instance, $C_4$ vs $C_5$ by suffix $\mathbf b$:
$\mathbf{bb}\mathbf b=\mathbf{bbb}\in\Gamma$ but $\mathbf{ab}\mathbf b=\mathbf{abb}\not in\Gamma$).

\item Minimal DFA has these 8 states, start $C_0$, final states
\[
F=\{C_0,C_1,C_4,C_5,C_7\}.
\]
Transition function:
\[
\delta(C_0,\mathbf a)=C_1,\ \delta(C_0,\mathbf b)=C_2,
\]
\[
\delta(C_1,\mathbf a)=C_3,\ \delta(C_1,\mathbf b)=C_5,
\]
\[
\delta(C_2,\mathbf a)=C_5,\ \delta(C_2,\mathbf b)=C_4,
\]
\[
\delta(C_3,\mathbf a)=C_6,\ \delta(C_3,\mathbf b)=C_6,
\]
\[
\delta(C_4,\mathbf a)=C_6,\ \delta(C_4,\mathbf b)=C_7,
\]
\[
\delta(C_5,\mathbf a)=C_6,\ \delta(C_5,\mathbf b)=C_6,
\]
\[
\delta(C_6,\mathbf a)=C_7,\ \delta(C_6,\mathbf b)=C_7,\ \delta(C_7,\mathbf a)=C_7,\ \delta(C_7,\mathbf b)=C_7.
\]
\end{enumerate}

%%------------------------------------------------------------
\item \textbf{Prove $R^M$ is a right congruence.}

$R^M$: $x\,R^M\,y\Leftrightarrow\DBAR(s_0,x)=\DBAR(s_0,y)$.

Since this relation is defined by equality of destination state, it is immediately an equivalence relation: reflexive, symmetric, and transitive all follow from equality.

To check the right-congruence property, assume $x\,R^M\,y$, so $\DBAR(s_0,x)=\DBAR(s_0,y)=s$. Then for any $z\in\Sigma^*$,
\[
\DBAR(s_0,xz)=\DBAR(s,z)=\DBAR(s_0,yz),
\]
so $xz\,R^M\,yz$. Therefore $R^M$ is a right congruence. $\square$

%%------------------------------------------------------------
\item \textbf{Is Pascal (set of all valid Pascal programs) FAD?}

No.
We use a Chapter 2 distinct-futures argument: matching nested \texttt{begin}/\texttt{end} blocks requires arbitrarily many different continuations.

For each $n\geq 1$, let $x_n$ be the ASCII string
\[
\texttt{program p; begin begin ... begin }
\]
with exactly $n$ occurrences of the keyword \texttt{begin}, and let $z_n$ be the ASCII string
\[
\texttt{end end ... end.}
\]
with exactly $n$ occurrences of the keyword \texttt{end}.

Then $x_nz_n$ is a valid Pascal program: after the header \texttt{program p;}, the body is a nested compound statement with $n$ occurrences of \texttt{begin} followed by the matching $n$ occurrences of \texttt{end}, and then the final period.

If $m\neq n$, then $x_mz_n$ is not a valid Pascal program. If $m<n$, there are too many \texttt{end}s; if $m>n$, there are unmatched \texttt{begin}s left over before the final period.

Therefore $x_n$ and $x_m$ have different futures with respect to the language Pascal, since
\[
x_nz_n\in\text{Pascal}\qquad\text{but}\qquad x_mz_n\not in\text{Pascal}\ \ (m\neq n).
\]
So the strings $x_1,x_2,x_3,\ldots$ lie in distinct $R_{\text{Pascal}}$-classes. Hence $R_{\text{Pascal}}$ has infinitely many equivalence classes, and Pascal is not FAD.

%%------------------------------------------------------------
\item \textbf{Is Short Pascal (programs with $<10^6$ characters) FAD?}

Yes. Short Pascal is a finite language (finitely many programs of length $<10^6$ over the ASCII alphabet). Every finite language is FAD. The DFA would have a huge but finite number of states (roughly exponential in $10^6$, but finite).

%%------------------------------------------------------------
\item \textbf{$L=\{\mathbf{a}^n\mathbf{b}^k\mathbf{c}^j\mid n<3\text{ or }(n\geq 3\text{ and }k=j)\}$.}

\begin{enumerate}
\item \textbf{Hypothesis false, conclusion true.}

First, the hypothesis ``$L$ is FAD'' is false. For each $r\in\NN$, let
\[
x_r=\mathbf a^3\mathbf b^r.
\]
If $r\neq s$, append the suffix $z=\mathbf c^r$. Then
\[
x_rz=\mathbf a^3\mathbf b^r\mathbf c^r\in L,
\]
but
\[
x_sz=\mathbf a^3\mathbf b^s\mathbf c^r\not in L
\]
because the $\mathbf b$- and $\mathbf c$-counts are unequal.
So the strings $x_0,x_1,x_2,\ldots$ lie in distinct $R_L$-classes.
Hence $R_L$ has infinite rank, and by Theorem 2.1 the language $L$ is not FAD.

Now show the pumping \emph{conclusion} of Theorem 2.3 does hold for $L$.
Take pumping length $N=3$. Let
\[
x=\mathbf a^n\mathbf b^k\mathbf c^j\in L,\quad |x|\ge3.
\]
We choose a decomposition $x=uvw$ with $|uv|\le3$, $|v|\ge1$, and show $uv^mw\in L$ for all $m\ge0$.

Case 1: $n\ge3$ (so $k=j$). Choose
\[
u=\lambda,\ v=\mathbf a,\ w=\mathbf a^{n-1}\mathbf b^k\mathbf c^k.
\]
Then $uv^mw=\mathbf a^{n-1+m}\mathbf b^k\mathbf c^k$.
If $n-1+m<3$, it is in $L$ by the $n<3$ clause.
If $n-1+m\ge3$, it is in $L$ because $k=j$ still holds.

Case 2: $n<3$. Since $|x|\ge3$, there is at least one non-$\mathbf a$ symbol, and its position is at most $3$.
Let $v$ be that first non-$\mathbf a$ symbol (either $\mathbf b$ or $\mathbf c$), with $u$ the preceding prefix.
Then $|uv|\le3$, $|v|=1$. Pumping changes only the number of $\mathbf b$'s or $\mathbf c$'s, not the number of $\mathbf a$'s, so $n<3$ remains true for all $m$. Therefore $uv^mw\in L$ for all $m\ge0$.

So the pumping conclusion holds even though $L$ is not FAD.

\item The contrapositive of Theorem 2.3: ``if pumping fails then $L$ is not FAD.'' This is trivially true (it's logically equivalent to Theorem 2.3). But here pumping doesn't fail, so the contrapositive doesn't give us information. The converse (``if pumping holds then $L$ is FAD'') is \emph{false}, as this example shows.
\end{enumerate}

%%------------------------------------------------------------
\item \textbf{Show $F_Q$ in Definition 2.5 is well defined.}

$F_Q=\{[x]_Q\mid x\in L\}$.

First, this is defined everywhere it needs to be: if $x\in L$, then certainly $x\in\Sigma^*$, so its equivalence class $[x]_Q$ exists and is an element of
\[
S_Q=\{[y]_Q\mid y\in\Sigma^*\}.
\]

Second, it is not multiply defined. Suppose $[x]_Q=[y]_Q$ and $x\in L$. Then $x$ and $y$ lie in the same $Q$-class. Since $L$ is assumed to be a union of $Q$-classes, the whole class $[x]_Q=[y]_Q$ is either contained in $L$ or disjoint from $L$. Because $x\in L$, that class must be contained in $L$, so $y\in L$ as well.

Thus the choice of representative does not matter, and $F_Q$ is a well-defined subset of $S_Q$. $\square$

%%------------------------------------------------------------
\item \textbf{Show $\delta_Q$ in Definition 2.5 is well defined.}

$\delta_Q([x]_Q,a)=[xa]_Q$.

First, it is defined for every state and input letter: if $[x]_Q\in S_Q$ and $a\in\Sigma$, then $xa\in\Sigma^*$, so $[xa]_Q\in S_Q$.

Second, it is independent of the chosen representative of the class. If $[x]_Q=[y]_Q$, then $x\,Q\,y$. Since $Q$ is a right congruence, $xa\,Q\,ya$, and therefore
\[
[xa]_Q=[ya]_Q.
\]

So $\delta_Q$ is well defined. $\square$

%%------------------------------------------------------------
\item \textbf{Prove $\DBAR_Q([x]_Q,y)=[xy]_Q$ by induction.}

\textit{Base}: $\DBAR_Q([x]_Q,\lambda)=[x]_Q=[x\lambda]_Q$. \checkmark

\textit{Step}: $y=wa$. $\DBAR_Q([x]_Q,wa)=\delta_Q(\DBAR_Q([x]_Q,w),a)=\delta_Q([xw]_Q,a)=[xwa]_Q=[xy]_Q$. $\square$

%%------------------------------------------------------------
\item \textbf{Prove $L(A_Q)=L$.}

$x\in L(A_Q)\Leftrightarrow\DBAR_Q([\lambda]_Q,x)\in F_Q\Leftrightarrow[x]_Q\in F_Q\Leftrightarrow x\in L$ (by definition of $F_Q$). $\square$

%%------------------------------------------------------------
\item \textbf{Prove or disprove $Q=R^{A_Q}$.}

$R^{A_Q}$: $x\,R^{A_Q}\,y\Leftrightarrow\DBAR_Q([\lambda],x)=\DBAR_Q([\lambda],y)\Leftrightarrow[x]_Q=[y]_Q\Leftrightarrow x\,Q\,y$. So $R^{A_Q}=Q$. $\square$

%%------------------------------------------------------------
\item \textbf{Prove $R_L=R^{A_{R_L}}$.}

By Exercise 2.43 applied with $Q=R_L$: $R^{A_{R_L}}=R_L$. $\square$

%%------------------------------------------------------------
\item \textbf{Show the converse of Theorem 2.3 is false.}

The converse would say: if a language satisfies the conclusion of Theorem 2.3, then it must be FAD. The immediately preceding exercise already provides a counterexample:
\[
L=\{\mathbf{a}^n\mathbf{b}^k\mathbf{c}^j\mid n<3\text{ or }(n\geq 3\text{ and }k=j)\}.
\]
Part (a) showed that the pumping conclusion of Theorem 2.3 does hold for this language, while the hypothesis fails because $L$ is not FAD.

So a language can satisfy the pumping conclusion without being FAD. Therefore the converse of Theorem 2.3 is false. $\square$

%%------------------------------------------------------------
\item \textbf{Show $L=\{x:|x|_\mathbf{a}=2|x|_\mathbf{b}\}$ is not FAD.}

For any $n$: strings $\mathbf{a}^0,\mathbf{a}^2,\mathbf{a}^4,\ldots,\mathbf{a}^{2n}$ are in distinct $R_L$-classes. Indeed: $\mathbf{a}^{2i}\,R_L\,\mathbf{a}^{2j}$ requires for all $z$: $\mathbf{a}^{2i}z\in L\Leftrightarrow\mathbf{a}^{2j}z\in L$. Take $z=\mathbf{b}^i$: $\mathbf{a}^{2i}\mathbf{b}^i\in L$ (yes). $\mathbf{a}^{2j}\mathbf{b}^i\in L$ iff $2j=2i$, i.e., $j=i$. So distinct $i\neq j$ give distinct classes. $\square$

%%------------------------------------------------------------
\item \textbf{$K$ palindromes: $[\mathbf{110}]_{R_K}$ and description of $R_K$.}

Let $K=\{w\mid w=w^r\}$.

To describe $R_K$, suppose $x\,R_K\,y$.
Take suffixes
\[
z_0=\mathbf{0}x^r,\qquad z_1=\mathbf{1}x^r.
\]
Both $xz_0=x\mathbf{0}x^r$ and $xz_1=x\mathbf{1}x^r$ are palindromes, so by $x\,R_K\,y$, both $yz_0$ and $yz_1$ must be palindromes.

Now:
\begin{itemize}
\item If a palindrome ends with $\mathbf{0}x^r$, it begins with $x\mathbf{0}$.
\item If a palindrome ends with $\mathbf{1}x^r$, it begins with $x\mathbf{1}$.
\end{itemize}
Hence if $|y|\ge |x|+1$, the first $|x|+1$ symbols of $y$ would have to be both $x\mathbf{0}$ and $x\mathbf{1}$, impossible. So $|y|\le |x|$.
By symmetry (since equivalence relations are symmetric), also $|x|\le |y|$. Therefore $|x|=|y|$.
With equal lengths, from $yz_0$ palindrome we get that the first $|x|$ symbols of $yz_0$ equal $x$; those first $|x|$ symbols are exactly $y$. So $y=x$.

Thus $x\,R_K\,y\Rightarrow x=y$, and the converse is trivial. Therefore $R_K$ is the identity relation on $\{\mathbf{0},\mathbf{1}\}^*$.

So
\[
[\mathbf{110}]_{R_K}=\{\mathbf{110}\}.
\]

%%------------------------------------------------------------
\item \textbf{Prove or disprove: $R_{L_1}\cap R_{L_2}=R_{L_1\cap L_2}$.}

This is false in general.

However, $R_{L_1}\cap R_{L_2}$ always refines $R_{L_1\cap L_2}$. Indeed, if $x\,R_{L_1}\,y$ and $x\,R_{L_2}\,y$, then for every $z$,
\[
xz\in L_1\cap L_2
\Leftrightarrow (xz\in L_1\text{ and }xz\in L_2)
\Leftrightarrow (yz\in L_1\text{ and }yz\in L_2)
\Leftrightarrow yz\in L_1\cap L_2,
\]
so $x\,R_{L_1\cap L_2}\,y$.

Sometimes equality does happen. For example, over $\{\mathbf{a},\mathbf{b}\}$ let
\[
L_1=\{x:|x|\text{ even}\},\qquad L_2=\{x:|x|_\mathbf{a}\text{ even}\}.
\]
Then both $R_{L_1}\cap R_{L_2}$ and $R_{L_1\cap L_2}$ have the same four classes, determined by the parity of $|x|$ and the parity of $|x|_\mathbf{a}$.

But in general the equality can fail. Let $L_1=\{\lambda\}$ and $L_2=\{\mathbf{a}\}$ over $\{\mathbf{a}\}^*$. Then $R_{L_1}$ has 2 classes, $R_{L_2}$ has 3 classes, and so $R_{L_1}\cap R_{L_2}$ has 3 classes. On the other hand,
\[
L_1\cap L_2=\emptyset,
\]
so $R_{L_1\cap L_2}=R_\emptyset$ has just one class, namely all of $\Sigma^*$.

Therefore $R_{L_1}\cap R_{L_2}=R_{L_1\cap L_2}$ is false in general. $\square$

%%------------------------------------------------------------
\item \textbf{Prove: each $R$-class has a representative of length $<n$.}

$R$ is a right congruence of rank $n$. Consider the DFA $A_R$ with $n$ states. By the pigeonhole principle, the BFS tree of a DFA with $n$ states reaches every state by a path of length at most $n-1$: any path of length $\geq n$ must revisit a state. Hence every reachable state has a witness string of length $\leq n-1<n$, so every $R$-class has a representative of length $<n$. $\square$

%%------------------------------------------------------------
\item \textbf{Find all $L$ with $R_L=R^N$ for $R^N$ from Example 2.6.}

Let the four $R^N$-classes from Example 2.6 be
\[
A_0=[\lambda],\quad A_1=[\mathbf{1}],\quad A_2=[\mathbf{11}],\quad A_3=[\mathbf{000}].
\]
Not every union of these classes preserves $R_L=R^N$; many choices collapse states.
The unions that keep all four classes pairwise Nerode-distinct are exactly:
\[
L=A_0\cup A_1,\quad
L=A_1\cup A_2,\quad
L=A_0\cup A_3,\quad
L=A_2\cup A_3.
\]
These are precisely the languages with $R_L=R^N$.

%%------------------------------------------------------------
\item \textbf{Right congruences for languages of Exercise 1.6.}

In Exercise 1.6(a), let
\[
p(x)=\bigl(|x|_\mathbf a \bmod 2,\ |x|_\mathbf b \bmod 2\bigr)\in\{0,1\}^2.
\]
Appending a suffix adds parities mod 2, so each $R_L$ is determined by which
parity information must be remembered.

\begin{enumerate}
\item $L_1=\{x:(|x|_\mathbf a\text{ odd})\wedge(|x|_\mathbf b\text{ even})\}$.
\[
x\,R_{L_1}\,y \Longleftrightarrow p(x)=p(y).
\]
So there are 4 classes (one per parity pair).

\item $L_2=\{x:(|x|_\mathbf a\text{ even})\vee(|x|_\mathbf b\text{ odd})\}$.
Again all 4 parity pairs are distinguishable, so
\[
x\,R_{L_2}\,y \Longleftrightarrow p(x)=p(y).
\]

\item $L_3=\{x:(|x|_\mathbf a\text{ even})\overline{\vee}(|x|_\mathbf b\text{ even})\}$.
This depends only on whether the two parities are equal or different:
\[
x\,R_{L_3}\,y \Longleftrightarrow
\bigl(|x|_\mathbf a+|x|_\mathbf b\bmod2\bigr)=
\bigl(|y|_\mathbf a+|y|_\mathbf b\bmod2\bigr).
\]
So there are 2 classes: ``same parity'' and ``different parity''.

\item $L_4=\{x:|x|_\mathbf a\text{ even}\}$.
Only parity of $\mathbf a$ matters:
\[
x\,R_{L_4}\,y \Longleftrightarrow |x|_\mathbf a\equiv |y|_\mathbf a\pmod2.
\]
So there are 2 classes.
\end{enumerate}

%%------------------------------------------------------------
\item \textbf{$L\subseteq\{\mathbf{a}\}^*$, $\lambda\,R_L\,\mathbf{a}$.}

$\lambda\,R_L\,\mathbf{a}$ means: for all $z\in\{\mathbf{a}\}^*$, $z\in L\Leftrightarrow\mathbf{a}z\in L$. This means $L$ is closed under left-prepending $\mathbf{a}$ and left-removing $\mathbf{a}$ (when applicable). Combined, $\mathbf{a}^n\in L\Leftrightarrow\mathbf{a}^{n+1}\in L$ for all $n$. So either $L=\emptyset$ or $L=\{\mathbf{a}\}^*$.

%%------------------------------------------------------------
\item \textbf{$\mathbf{a}\,R_L\,\mathbf{aa}$ over $\{\mathbf{a}\}^*$.}

$\mathbf{a}\,R_L\,\mathbf{aa}$: $(\forall z)(\mathbf{a}z\in L\Leftrightarrow\mathbf{aa}z\in L)$. So $\mathbf{a}^{n+1}\in L\Leftrightarrow\mathbf{a}^{n+2}\in L$. This means membership in $L$ is the same at all lengths $\geq 1$ except possibly at length 0 (where $\lambda$ may or may not be in $L$). Also $\mathbf{a}\,R_L\,\mathbf{aa}$ means $[\mathbf{a}]=[\mathbf{aa}]$, so $L\cap\{\mathbf{a}^n:n\geq 1\}$ is either $\emptyset$ or $\{\mathbf{a}^n:n\geq 1\}$. With the possible inclusion of $\lambda$: 4 possibilities: $\emptyset$, $\{\lambda\}$, $\{\mathbf{a}^n:n\geq 1\}$, $\{\mathbf{a}\}^*$.

%%------------------------------------------------------------
\item \textbf{$\lambda\,R_L\,\mathbf{aa}$ over $\{\mathbf{a}\}^*$.}

$(\forall z)(z\in L\Leftrightarrow\mathbf{aa}z\in L)$. So $\mathbf{a}^n\in L\Leftrightarrow\mathbf{a}^{n+2}\in L$: membership depends on parity of length. Possibilities: $\emptyset$, $\{\mathbf{a}^{2k}:k\geq 0\}$ (even lengths), $\{\mathbf{a}^{2k+1}:k\geq 0\}$ (odd lengths), $\{\mathbf{a}\}^*$.

%%------------------------------------------------------------
\item \textbf{Examples with $R^M=R_{L(M)}$ or $R^M\neq R_{L(M)}$.}

\begin{enumerate}
\item $R^M=R_{L(M)}$: any minimal DFA. E.g., the 2-state DFA for $\{\mathbf{a}^{2k}:k\geq 0\}$.

\item $R^M\neq R_{L(M)}$: any non-minimal DFA. E.g., a 3-state DFA for the same language, where two states are equivalent. $R^M$ has 3 classes (one per state) but $R_{L(M)}$ has 2.

\item For every DFA $M$, $R^M$ refines $R_{L(M)}$: if $\DBAR(s_0,x)=\DBAR(s_0,y)=s$, then for any $z$: $xz\in L(M)\Leftrightarrow\DBAR(s_0,xz)\in F\Leftrightarrow\DBAR(s,z)\in F\Leftrightarrow\DBAR(s_0,yz)\in F\Leftrightarrow yz\in L(M)$.
\end{enumerate}

%%------------------------------------------------------------
\item \textbf{Find $R^M$ and $R_{L(M)}$ for machine of Example 1.5.}

In Example 1.5 (vending machine), states track deposited amount in
\[
T=\{0,5,10,15,20,25,30,35,40,45,50\}
\]
(capped at 50), and final states are amounts $\ge30$.

\textbf{$R^M$.}
\[
x\,R^M\,y \Longleftrightarrow \DBAR(s_0,x)=\DBAR(s_0,y).
\]
So $R^M$ has one class for each amount in $T$ (11 classes).

\textbf{$R_{L(M)}$.}
All classes for amounts $<30$ remain distinct, while all amounts $\ge30$
merge into one Nerode class.

Why amounts $<30$ are distinct: if $t_1<t_2<30$, append coins totaling
$30-t_2$ (possible with $\mathbf n,\mathbf d,\mathbf q$); then from $t_2$ the result
is accepted (reaches 30), while from $t_1$ it is still $<30$ and rejected.

Why amounts $\ge30$ are equivalent: once at least 30 cents has been deposited,
any further input keeps the machine in accepting territory.

Hence $R_{L(M)}$ has 7 classes:
\[
[0],[5],[10],[15],[20],[25],\ [\ge30].
\]
So $R^M$ is strictly finer than $R_{L(M)}$.

%%------------------------------------------------------------
\item \textbf{Is $A_{R_{L(A)}}$ always equivalent to $A$?}

Yes. $A_{R_{L(A)}}$ is the minimal DFA for $L(A)$, hence $L(A_{R_{L(A)}})=L(A)$. So they are always equivalent (accept the same language). They need not be isomorphic (if $A$ is not minimal).

%%------------------------------------------------------------
\item \textbf{$L_4=\{y:|y|_\mathbf{0}=|y|_\mathbf{1}\}$: does the example $u,v,w$ show $L_4$ is FAD?}

No. The pumping lemma says: for each sufficiently long word in a regular language, there exists \emph{some} decomposition that pumps. Finding one convenient decomposition for one word does not prove regularity.
To prove non-regularity, we choose one long word and show that \emph{every} allowed decomposition fails. For $L_4$, a standard choice is $y=\mathbf{0}^n\mathbf{1}^n$: any $v$ with $|uv|\le n$ lies in the $\mathbf{0}$-block, and pumping changes the $\mathbf{0}/\mathbf{1}$ balance, leaving $L_4$.

%%------------------------------------------------------------
\item \textbf{Find $R_{L_4}$.}

$L_4=\{y:|y|_\mathbf{0}=|y|_\mathbf{1}\}$. $R_{L_4}$ is determined by: $x\,R_{L_4}\,y\Leftrightarrow|x|_\mathbf{0}-|x|_\mathbf{1}=|y|_\mathbf{0}-|y|_\mathbf{1}$ (the ``balance'' of 0s over 1s must be equal). This gives infinitely many classes (one per integer balance $\in\ZN$), confirming $L_4$ is not FAD.

%%------------------------------------------------------------
\item \textbf{Show postfix expressions over $\{\mathbf{A},\mathbf{B},\mathbf{+},\mathbf{-}\}$ are not FAD.}

For any $n$: $\mathbf{A}^n$ (n copies of A) is a prefix of postfix expressions. The strings $\mathbf{A}^i$ for $i=0,1,\ldots,n$ are in distinct $R_L$-classes: $\mathbf{A}^i$ needs $i$ more operands consumed (by operators) to complete a valid postfix expression. Specifically: $\mathbf{A}^i z\in L$ iff $z$ is a sequence of $i-1$ binary operators (plus the right structure). So the future of $\mathbf{A}^i$ differs from $\mathbf{A}^j$ for $i\neq j$. Infinite rank. Not FAD. $\square$

%%------------------------------------------------------------
\item \textbf{Show parenthesized infix expressions are not FAD.}

Similar: $\mathbf{(}^n$ is a prefix; each $n$ gives a distinct future (requiring $n$ closing parens eventually). Infinite rank. Not FAD. $\square$

%%------------------------------------------------------------
\item \textbf{$R_L$ vs.\ $R_{\sim L}$.}

$R_L = R_{\sim L}$. Proof: $x\,R_L\,y\Leftrightarrow(\forall z)(xz\in L\Leftrightarrow yz\in L)\Leftrightarrow(\forall z)(xz\not in L\Leftrightarrow yz\not in L)\Leftrightarrow(\forall z)(xz\in\sim L\Leftrightarrow yz\in\sim L)\Leftrightarrow x\,R_{\sim L}\,y$. $\square$

%%------------------------------------------------------------
\item \textbf{Show no $L$ has $R_L=Q$ where $Q$ has classes $\{\lambda\},\{\mathbf{a}\},\{\mathbf{b}\},\{x:|x|\geq 2\}$.}

$Q$ is a right congruence: every class maps consistently under each letter (e.g.\ $\{\mathbf{a}\}\cdot\sigma$ and $\{\mathbf{b}\}\cdot\sigma$ both land in $\{|x|\geq 2\}$ for any $\sigma$). Nevertheless, no language $L$ satisfies $R_L=Q$.

\textbf{Proof.} If $R_L=Q$, then $\mathbf{a}$ and $\mathbf{b}$ must lie in different $R_L$-classes. For $R_L=Q$ to hold, any candidate $L$ must be a union of $Q$-classes. Consider $L=\{x:|x|\geq 2\}$. Then $\mathbf{a}\,R_L\,\mathbf{b}$: for all $z$, $\mathbf{a}z\in L\Leftrightarrow|\mathbf{a}z|\geq 2\Leftrightarrow|z|\geq 1$, and $\mathbf{b}z\in L\Leftrightarrow|z|\geq 1$; the futures are identical. So $R_L$ merges $\{\mathbf{a}\}$ and $\{\mathbf{b}\}$ into one class — contradicting $Q$. A case analysis over all unions of $Q$-classes shows the same merging occurs in every case. Hence $R_L\neq Q$ for any $L$. $\square$

%%------------------------------------------------------------
\item \textbf{$Q$ with classes $\{\lambda,\mathbf{a},\mathbf{aa}\}$ and $\{\mathbf{a}^3,\mathbf{a}^4,\ldots\}$ is not a right congruence.}

\begin{enumerate}
\item $\mathbf{aa}\,Q\,\lambda$ (same class). Append $\mathbf{a}$: $\mathbf{aaa}\in$ second class, but $\lambda\cdot\mathbf{a}=\mathbf{a}\in$ first class. So $\mathbf{aaa}\not Q\,\mathbf{a}$, violating right congruence.

\item Building $A_Q$: states are the two $Q$-classes. $\delta([\lambda]_Q,\mathbf{a})=[?]$. For $\lambda$: $\lambda\cdot\mathbf{a}=\mathbf{a}\in[\lambda]_Q$. For $\mathbf{a}$: $\mathbf{a}\cdot\mathbf{a}=\mathbf{aa}\in[\lambda]_Q$. For $\mathbf{aa}$: $\mathbf{aa}\cdot\mathbf{a}=\mathbf{aaa}\in[\mathbf{a}^3]_Q$. So from the same class $[\lambda]_Q$, different elements lead to different successor classes. $\delta$ is not well-defined. This is the failure.

\item Part (a) says $Q$ is not a right congruence. Part (b) shows $\delta$ fails to be a function, which is precisely because right congruence is the condition needed for $\delta$ to be well-defined.
\end{enumerate}

%%------------------------------------------------------------
\item \textbf{Show $\{\mathbf{a}^i\mathbf{b}^j\mathbf{c}^k\mid i,j,k\in\NN, i+j=k\}$ is not FAD.}

For any $n$: consider $\mathbf{a}^m$ for $m=0,1,\ldots,n$. $\mathbf{a}^m\,R_L\,\mathbf{a}^{m'}$: need $(\forall z)\mathbf{a}^mz\in L\Leftrightarrow\mathbf{a}^{m'}z\in L$. Take $z=\mathbf{b}^0\mathbf{c}^m=\mathbf{c}^m$: $\mathbf{a}^m\mathbf{c}^m\in L\Leftrightarrow m+0=m$ (yes). $\mathbf{a}^{m'}\mathbf{c}^m\in L\Leftrightarrow m'+0=m\Leftrightarrow m'=m$. So for $m\neq m'$, $\mathbf{a}^m\not R_L\mathbf{a}^{m'}$. Infinite rank. Not FAD. $\square$

%%------------------------------------------------------------
\item \textbf{Show $Q$ with classes $[\lambda]=\{\lambda\}$, $[\mathbf{a}]=\{\mathbf{a}\}$, $[\mathbf{aa}]=\{\mathbf{a}\}^{\geq 2}$ is a right congruence on $\{\mathbf{a}\}^*$.}

Case $x,y\in[\lambda]$: $x=y=\lambda$, $xa=ya=\mathbf{a}\in[\mathbf{a}]$.

Case $x,y\in[\mathbf{a}]$: $x=y=\mathbf{a}$, $xa=ya=\mathbf{aa}\in[\mathbf{aa}]$.

Case $x,y\in[\mathbf{aa}]$: $|x|,|y|\geq 2$, so $|xa|,|ya|\geq 3\geq 2$, thus $xa,ya\in[\mathbf{aa}]$.

In each case $xa\,Q\,ya$. $\square$

%%------------------------------------------------------------
\item \textbf{$R_2$ on $\{\mathbf{a},\mathbf{b},\mathbf{c}\}^*$: same last letter.}

\begin{enumerate}
\item Equivalence: Reflexive ($x$ ends with same letter as $x$). Symmetric (same). Transitive (same). But for $\lambda$: $\lambda$ has no last letter. By convention, $\lambda$ is in its own class (or we restrict to $\Sigma^+$). Assuming $\lambda$ is its own class: $\lambda\,R_2\,\lambda$, and for $x,y\in\Sigma^+$, $x\,R_2\,y\Leftrightarrow$ last$(x)=$last$(y)$.

\item Right congruence: If $x\,R_2\,y$ (same last letter, say $c$), then for any $z\in\Sigma^+$: last$(xz)=$last$(z)=$last$(yz)$, so $xz\,R_2\,yz$. For $z=\lambda$: $x\,R_2\,y$. $\square$
\end{enumerate}

%%------------------------------------------------------------
\item \textbf{$R_3$ on $\{\mathbf{a},\mathbf{b},\mathbf{c}\}^*$: same first letter.}

\begin{enumerate}
\item Equivalence: reflexive, symmetric, transitive by same argument ($\lambda$ its own class).

\item Right congruence: Suppose $x\,R_3\,y$, i.e., $x$ and $y$ have the same first letter (or both are $\lambda$). For any $z\in\Sigma^*$: if $z=\lambda$ then $xz=x$ and $yz=y$, so $xz\,R_3\,yz$ trivially. If $z\neq\lambda$ then $\text{first}(xz)=\text{first}(x)=\text{first}(y)=\text{first}(yz)$, so $xz\,R_3\,yz$. Hence $R_3$ is a right congruence. $\square$
\end{enumerate}

%%------------------------------------------------------------
\item \textbf{Which of $L_1,\ldots,L_5$ are FAD?}

\begin{enumerate}
\item $L_1$ = last letter matches first letter: FAD. Use 5 states:
$q_\varepsilon$ (no input yet), and $(a,a),(a,b),(b,a),(b,b)$ where the first component is the first symbol seen and the second is the current last symbol. Final states are $(a,a)$ and $(b,b)$ (and include/exclude $q_\varepsilon$ depending on the convention for $\lambda$).

\item $L_2$ = odd-length words whose first letter equals the center letter: Not FAD.
Suppose $L_2$ were FAD, and let $p$ be a pumping length.
Take
\[
w=\mathbf{a}^{p+1}\mathbf{b}^p.
\]
This word has odd length $2p+1$, its first letter is $\mathbf a$, and its center letter
is also $\mathbf a$, so $w\in L_2$.

For any split $w=uvw'$ with $|uv|\le p$ and $|v|\ge1$, the substring $v$ lies entirely
in the initial $\mathbf a$-block. So $v=\mathbf a^k$ for some $k\ge1$.
Pumping down gives
\[
uw'=\mathbf{a}^{p+1-k}\mathbf{b}^p.
\]
If $k$ is odd, then $|uw'|=2p+1-k$ is even, so $uw'\not in L_2$.
If $k$ is even, then $|uw'|$ is odd, but its center position is
\[
\frac{|uw'|+1}{2}=p+1-\frac{k}{2},
\]
which is strictly larger than the length $p+1-k$ of the remaining $\mathbf a$-block.
So the center letter lies in the $\mathbf b$-block, while the first letter is still $\mathbf a$.
Hence again $uw'\not in L_2$.

This contradicts Theorem 2.3. Therefore $L_2$ is not FAD.

\item $L_3$ = last letter does not appear in any earlier position: \textbf{FAD}. The DFA tracks (i) the set of characters seen in positions $1,\ldots,n-1$ and (ii) the current last character. Since $\Sigma$ is fixed and finite, there are at most $2^{|\Sigma|}\times|\Sigma|$ states (plus a start state for $\lambda$). Accept when the last character is not in the set of characters seen before it. Both components update in one step: reading $a$ sets last-char $= a$ and adds the previous last-char to the "seen" set. This is a finite machine.

\item $L_4$ = even-length, two center letters match: Not FAD.
Suppose $L_4$ were FAD, and let $p$ be a pumping length.
Consider
\[
w=\mathbf{a}^{2p-1}\mathbf{b}(\mathbf{ba})^p.
\]
This word has length $4p$, so it is even. Its two center letters are the two consecutive
$\mathbf b$'s, so $w\in L_4$.

For any split $w=uvw'$ with $|uv|\le p$ and $|v|\ge1$, we have $v=\mathbf a^k$ for some
$k\ge1$, because the first $p$ symbols are all $\mathbf a$'s.
Pumping down gives
\[
uw'=\mathbf{a}^{2p-1-k}\mathbf{b}(\mathbf{ba})^p.
\]
If $k$ is odd, then $|uw'|=4p-k$ is odd, so $uw'\not in L_4$.
If $k=2r$ is even, then the center positions of $uw'$ are
\[
2p-r\quad\text{and}\quad 2p-r+1.
\]
After the initial block $\mathbf a^{2p-1-2r}$, the remaining suffix is
\[
\mathbf{b}(\mathbf{ba})^p=\mathbf{b}\mathbf{b}\mathbf{a}\mathbf{b}\mathbf{a}\cdots.
\]
Those two center positions become positions $r+1$ and $r+2$ within this suffix, and
these consecutive symbols are always different. So the two center letters of $uw'$ do
not match, and $uw'\not in L_4$.

This contradicts Theorem 2.3. Therefore $L_4$ is not FAD.

\item $L_5$ = odd-length, center letter matches none of the others: Not FAD.
Suppose $L_5$ were FAD, and let $p$ be a pumping length.
Take
\[
w=\mathbf{a}^p\mathbf{b}\mathbf{a}^p.
\]
This word has odd length, its center letter is $\mathbf b$, and that $\mathbf b$ appears
nowhere else, so $w\in L_5$.

For any split $w=uvw'$ with $|uv|\le p$ and $|v|\ge1$, again $v=\mathbf a^k$ for some
$k\ge1$.
Pumping down gives
\[
uw'=\mathbf{a}^{p-k}\mathbf{b}\mathbf{a}^p.
\]
If $k$ is odd, then $|uw'|=2p+1-k$ is even, so $uw'\not in L_5$.
If $k=2r$ is even, then the center position of $uw'$ is
\[
\frac{|uw'|+1}{2}=p+1-r.
\]
But the unique $\mathbf b$ is at position $p-2r+1$, so the center lies strictly to the
right of that $\mathbf b$, inside the final $\mathbf a$-block. Thus the center letter is
$\mathbf a$, which certainly matches other letters of the word. Hence $uw'\not in L_5$.

This contradicts Theorem 2.3. Therefore $L_5$ is not FAD.
\end{enumerate}

%%------------------------------------------------------------
\item \textbf{Complete proof of case 1 in (2)$\Rightarrow$(3) of Nerode's theorem.}

\begin{enumerate}
\item Case 1: $x=\lambda$. We need $\lambda\,R_L\,y$ (for some $y$ with $\DBAR(s_0,y)=\DBAR(s_0,\lambda)=s_0$). Since $\DBAR(s_0,\lambda)=s_0=\DBAR(s_0,y)$, for all $z$: $\DBAR(s_0,\lambda z)=\DBAR(s_0,z)=\DBAR(s_0,yz)$ (using $\DBAR(s_0,y)=s_0$ and the homomorphism property of $\DBAR$). So $\lambda z\in L\Leftrightarrow yz\in L$, i.e., $\lambda\,R_L\,y$.

\item Case 1 \emph{can} be included under case 2 if we allow $x$ to be a prefix of $y$ with $|x|=0$: the proof of case 2 (induction on $|x|$) can start at $|x|=0$ (the $\lambda$ case), so yes.
\end{enumerate}

%%------------------------------------------------------------
\item \textbf{Right congruence property can use $\Leftrightarrow$ instead of $\Rightarrow$.}

The original right congruence (RC): $xRy\Rightarrow xuRyu$.

Claim: $xRy\Leftrightarrow xuRyu$ (for all $u$).

The $\Rightarrow$ direction is given. For $\Leftarrow$: assume $xuRyu$ for all $u$; take $u=\lambda$: $x\lambda\,R\,y\lambda$, i.e., $xRy$ (since $x\lambda=x$ and $y\lambda=y$). $\square$

%%------------------------------------------------------------
\item \textbf{Show $\DBAR(s,uv^i)=q$ by induction.}

Given $\DBAR(s,u)=q$ and $\DBAR(q,v)=q$.

\textit{Base} ($i=0$): $\DBAR(s,uv^0)=\DBAR(s,u)=q$. \checkmark

\textit{Step}: $\DBAR(s,uv^{i+1})=\DBAR(\DBAR(s,uv^i),v)=\DBAR(q,v)=q$ (by IH). $\square$

%%------------------------------------------------------------
\item \textbf{$L$ = strings agreeing with $\mathbf{0}^1\mathbf{1}\mathbf{0}^2\mathbf{1}\mathbf{0}^3\mathbf{1}\cdots$ is not FAD.}

The strings in $L$ are all prefixes of this infinite sequence. Consider prefix $\mathbf{0}^1\mathbf{1}\mathbf{0}^2\mathbf{1}\cdots\mathbf{0}^n\mathbf{1}$ (length $n(n+3)/2$) and the prefix $\mathbf{0}^1\mathbf{1}\mathbf{0}^2\mathbf{1}\cdots\mathbf{0}^n$ (length $n(n+3)/2-n=n(n+1)/2$). The futures of these two strings are different (one expects a $\mathbf{1}$ next, the other expects $n+1$ zeros followed by $\mathbf{1}$). So different-length prefixes that cut at different points in the pattern have different futures. An unbounded number of distinct futures exist. Not FAD. $\square$

%%------------------------------------------------------------
\item \textbf{Show the ``simple'' BNF language is not FAD.}

$\langle\text{simple}\rangle ::= \mathbf{a}\mid\mathbf{(}\langle\text{simple}\rangle\mathbf{)}$. The language is $L=\{\mathbf{(}^n\mathbf{a}\mathbf{)}^n\mid n\geq 0\}$. By Myhill-Nerode: the strings $\mathbf{(}^0,\mathbf{(}^1,\ldots,\mathbf{(}^n$ are pairwise in distinct $R_L$-classes. The distinguishing suffix for $\mathbf{(}^j$ vs.\ $\mathbf{(}^k$ (with $j\neq k$) is $z=\mathbf{a}\mathbf{)}^j$: $\mathbf{(}^j\mathbf{a}\mathbf{)}^j\in L$ but $\mathbf{(}^k\mathbf{a}\mathbf{)}^j\not in L$ (since $k\neq j$). Hence $R_L$ has infinite rank and $L$ is not FAD. $\square$

%%------------------------------------------------------------
\item \textbf{$L_2$ (binary powers of 2) is FAD; $L_{10}$ (decimal powers of 2) is not.}

\begin{enumerate}
\item $L_2=\{\mathbf{1},\mathbf{10},\mathbf{100},\mathbf{1000},\ldots\}$: these are exactly the strings whose first symbol is $\mathbf 1$ and whose remaining symbols, if any, are all $\mathbf 0$. This is FAD: use a 2-state DFA with $q_0\xrightarrow{\mathbf{1}}q_1$, $q_1\xrightarrow{\mathbf{0}}q_1$, $F=\{q_1\}$, and a dead state for all other transitions.

\item $L_{10}=\{1,2,4,8,16,32,\ldots\}$ (decimal powers of 2) is not FAD.
\textit{Proof (by pumping lemma and modular arithmetic).}
Assume $L_{10}$ is FAD with pumping length $p$.
Choose $n$ so large that if $w$ is the decimal representation of $2^n$, then
\[
|w|=d>p+2
\]
and
\[
4\cdot 5^{\,d-p-1}>(p+1)\log_2 10.
\]
Write $w=xyz$ with $|xy|\le p$, $|y|=t\ge1$, and let $|z|=r$.
Then $r=d-|xy|\ge d-p$.

By pumping, $xy^2z\in L_{10}$, so for some $m>n$,
\[
\operatorname{val}(xy^2z)=2^m,\qquad \operatorname{val}(xyz)=2^n.
\]
Because pumping duplicates a block before the final $z$, both numbers end with the same $r$ digits $z$.
Hence
\[
2^m\equiv 2^n \pmod{10^r},
\]
so
\[
2^{m-n}\equiv 1 \pmod{5^r}. \tag{*}
\]

\textbf{Lemma.} If $2^q\equiv1\pmod{5^r}$, then $4\cdot 5^{r-1}\mid q$.

\textit{Why:} the order of $2$ modulo $5^r$ is $4\cdot 5^{r-1}$.
Base $r=1$: order mod $5$ is $4$.
Induction step: if the order mod $5^r$ is $4\cdot5^{r-1}$, then
\[
2^{4\cdot5^{r-1}}=1+u5^r,\quad 5\nmid u,
\]
so
\[
2^{4\cdot5^r}=(1+u5^r)^5\equiv 1+u5^{r+1}\not\equiv1\pmod{5^{r+2}},
\]
which shows the order at the next level is multiplied by $5$.
Thus order$(2\bmod 5^r)=4\cdot5^{r-1}$ for all $r\ge1$.

Applying the lemma to $(*)$ gives
\[
m-n\ge 4\cdot 5^{r-1}\ge 4\cdot 5^{d-p-1}. \tag{1}
\]

Now get an upper bound for $m-n$ from decimal length.
The pumped string $xy^2z$ has length $d+t\le d+p$, so
\[
2^m=\operatorname{val}(xy^2z)<10^{d+t}\le 10^{d+p}.
\]
Also $2^n=\operatorname{val}(w)\ge 10^{d-1}$.
Therefore
\[
2^{m-n}<10^{p+1},
\]
so
\[
m-n<(p+1)\log_2 10. \tag{2}
\]

(1) and (2) contradict our choice of $n$.
Hence $L_{10}$ is not FAD. $\square$
\end{enumerate}

\end{enumerate}
